%======================================================================%
% UNIT: strings — Vortex Strings, Heterotic Correspondence, and Conclusions
%======================================================================%
\section{Vortex strings, heterotic correspondence, and conclusions}
\label{sec:strings:main}

We close the construction by collecting the four threads of the broken
phase --- topological strings, their worldsheet algebra, the heterotic
moduli match, and the Standard-Model descent --- into a single,
systematic classification of the principal results. Throughout we
keep the signature $(-,+,+,+)$, write $[f]=\text{mass}$ for the decay
constant, and call the $32$ coset Goldstone directions the \emph{clock
fields}. The vacuum manifold is the rank-$1$ primitive-idempotent orbit
\begin{equation}
  \mathcal{M} \;=\; E_{6}/(\mathrm{Spin}(10)\times U(1)) \;=\; \mathrm{EIII},
  \qquad \dim_{\mathbb{R}}\mathcal{M}=32,
  \label{eq:strings:vacuum}
\end{equation}
with positive-definite coset metric $K|_{\mathfrak m}$ and tangent module
$\mathfrak m = \mathbf{16}_{-3}\oplus\overline{\mathbf{16}}_{+3}$
(\autoref{sec:coset:main}). The physical spectrum is $28$ massive clock
scalars of a single common mass $m^{2}=\mu^{2}$ plus $4$ gauge/soldering
modes eaten by diffeomorphism and local Lorentz invariance
(\autoref{sec:gravity:main}).

%----------------------------------------------------------------------%
\subsection{The topological source of the strings}
\label{sec:strings:topology}

The naive expectation is that the broken phase supports cosmic strings
classified by $\pi_{1}(\mathcal{M})$. This must be stated with care.

\begin{theorem}[EIII is simply connected]
\label{thm:strings:pi1}
The vacuum manifold \eqref{eq:strings:vacuum} is a compact Hermitian
symmetric space (the complex Cayley plane). Hence
\[
   \pi_{1}(\mathrm{EIII})=0 .
\]
\end{theorem}

\begin{proof}
$\mathrm{EIII}=E_{6}/(\mathrm{Spin}(10)\times U(1))$ is an irreducible
compact Hermitian symmetric space of type $\mathrm{E\,III}$
(\'E.~Cartan's classification). The total space $E_{6}$ (compact, simply
connected) fibres over $\mathcal{M}$ with connected fibre
$\mathrm{Spin}(10)\times U(1)$, which is itself connected; the long exact
homotopy sequence
$\pi_{1}(E_{6})\to\pi_{1}(\mathcal{M})\to\pi_{0}(\mathrm{Spin}(10)\times
U(1))$ has both flanking groups trivial, so $\pi_{1}(\mathcal{M})=0$.
Equivalently, all compact irreducible Hermitian symmetric spaces are
simply connected.
\end{proof}

\begin{remark}[Topological consistency]
\label{rem:strings:guard}
\autoref{thm:strings:pi1} \emph{forbids} the assertion
$\pi_{1}(\mathcal{M})=\mathbb{Z}$ \emph{of EIII itself}: there are no
$\sigma$-model strings winding the bare coset target. Any source claiming
$\pi_{1}(\mathrm{EIII})=\mathbb{Z}$ is in error. The correct topological
seat of the strings is the \emph{gauged} configuration space, identified
next.
\end{remark}

The strings instead arise once the composite $\mathrm{Spin}(10)\times
U(1)$ connection (\autoref{sec:gauge:main}) is switched on. Gauging the
$U(1)\subset H$ that acts on the clock fields replaces the target by the
configuration space of the gauged $\sigma$-model. Concretely, the
physical order parameter for vortices is the $U(1)$-charged combination
of clock fields whose phase lives in
\[
   \mathcal{C}_{U(1)} \;\simeq\; \mathcal{M}\,/\!/\,U(1),
   \qquad
   \pi_{1}\bigl(\mathcal{C}_{U(1)}\bigr)\;\supseteq\;\mathbb{Z},
   \label{eq:strings:config}
\]
the winding being the standard Abelian-Higgs flux of the gauged $U(1)$
(the same $U(1)$ whose Stueckelberg/Green--Schwarz coupling appears in
\autoref{sec:anomaly:main}). Equivalently, on the broken vacuum the
$U(1)$ factor of $H=\mathrm{Spin}(10)\times U(1)$ admits a magnetic flux
tube: $\pi_{1}(U(1))=\mathbb{Z}$ labels the quantised flux, and the clock
field winds the eaten phase around the core.

\begin{postulate}[Topological string sector]
\label{post:strings:sector}
The relevant first homotopy is that of the $U(1)$-gauged configuration
space \eqref{eq:strings:config}, $\pi_{1}(\mathcal{C}_{U(1)})\supseteq
\mathbb{Z}$, \emph{not} that of the bare coset. The vortex strings are
flux tubes of the composite $U(1)$ (cf.\ Axiom~I and the gauging
prescription, \autoref{sec:gauge:main}); we treat the string content as a
\emph{correspondence} with Abelian-Higgs/heterotic vortices rather than a
$\pi_{1}(\mathcal{M})$ theorem.
\end{postulate}

%----------------------------------------------------------------------%
\subsection{BPS vortex strings of finite tension}
\label{sec:strings:bps}

Choose cylindrical coordinates $(r,\theta,z,t)$ on the emergent
Lorentzian manifold $(X,\eta)$ of \autoref{sec:soldering:main}. A static
configuration of the $U(1)$-gauged clock-field sector with winding
$n\in\mathbb{Z}$ obeys
\begin{equation}
  \phi(r,\theta+2\pi)=e^{\,n\,T_{U(1)}\,2\pi}\,\phi(r,\theta),
  \qquad
  \phi(0)=X_{0}\in\mathcal{M},\quad
  D_{r}\phi\xrightarrow[r\to\infty]{}0,
  \label{eq:strings:ansatz}
\end{equation}
with $T_{U(1)}$ the broken $U(1)$ generator and $X_{0}$ a rank-$1$
idempotent representative of the vacuum.

\begin{theorem}[Finite BPS tension]
\label{thm:strings:tension}
For each $n\in\mathbb{Z}$ the Euler--Lagrange equations of the gauged
clock-field action admit a static, axially symmetric, finite-tension
solution of the ansatz type \eqref{eq:strings:ansatz}. Its tension
saturates a Bogomol'nyi bound,
\begin{equation}
  \mathcal{T}_{n}
   = 2\pi f^{2}\,|n|\,
     \int_{0}^{\infty}\!dr\;
     r\,K_{\alpha\beta}(\phi)\,
     D_{r}\phi^{\alpha}\,D_{r}\phi^{\beta}
   \;=\; 2\pi f^{2}\,|n|\,\Omega,
  \qquad 0<\mathcal{T}_{n}<\infty,
  \label{eq:strings:tension}
\end{equation}
where $\Omega$ is the $U(1)$ flux quantum through the core and
$K_{\alpha\beta}$ the positive-definite coset metric.
\end{theorem}

\begin{proof}[Proof sketch]
Completing the static energy functional into a sum of squares plus a
topological flux term gives a first-order Bogomol'nyi system whose
solutions minimise the energy in each winding class; the
positive-definiteness of $K|_{\mathfrak m}$ (\autoref{sec:coset:main})
guarantees a strictly positive integrand, and the exponential approach
$D_{r}\phi\to0$ (set by the common mass $m^{2}=\mu^{2}$ of the $28$
massive clock fields) makes the integral converge. The detailed
moduli-matrix construction is deferred to the worldsheet appendix,
\autoref{app:ws:main}.
\end{proof}

Linear fluctuations about the BPS background factorise into right-moving
Nambu--Goto modes (transverse string position) and left-moving internal
currents (the eaten $\mathrm{Spin}(10)\times U(1)$ phases of the core).

%----------------------------------------------------------------------%
\subsection{Worldsheet algebra and the heterotic correspondence}
\label{sec:strings:worldsheet}

\begin{theorem}[Left-moving affine algebra]
\label{thm:strings:affine}
Quantisation of the gauge zero modes trapped on the vortex core yields a
left-moving affine current algebra at level one,
\begin{equation}
  \widehat{\mathfrak{so}}(10)_{1}\;\oplus\;\widehat{\mathfrak u}(1)_{1}.
  \label{eq:strings:affine}
\end{equation}
\end{theorem}

\begin{proof}[Proof sketch]
The Green--Schwarz term $S_{\mathrm{GS}}=\int B\wedge
\mathrm{Tr}(F\wedge F)/8\pi^{2}$ (\autoref{sec:anomaly:main}) reduces on
the two-dimensional string core $\Sigma\simeq\mathbb{R}^{1,1}$ to a
chiral WZW action whose level equals the mixed
$U(1)$--$\mathrm{Spin}(10)^{2}$ anomaly coefficient computed there;
matching that coefficient fixes $k=k'=1$. The Sugawara stress tensor of
\eqref{eq:strings:affine} then governs the left-movers. Full details
appear in \autoref{app:ws:main}.
\end{proof}

\begin{theorem}[Sugawara central charge]
\label{thm:strings:cl}
The Sugawara central charge of \eqref{eq:strings:affine} is
\begin{equation}
  c_{L}
   = \frac{\dim\mathrm{SO}(10)\cdot k}{k+h^{\vee}_{\mathrm{SO}(10)}}
     + 1
   = \frac{45\cdot 1}{1+8} + 1
   = 5 + 1
   = \boxed{16},
  \label{eq:strings:cL}
\end{equation}
using $\dim\mathrm{SO}(10)=45$, dual Coxeter number
$h^{\vee}_{\mathrm{SO}(10)}=8$, and $c[\widehat{\mathfrak u}(1)_{1}]=1$.
\end{theorem}

\begin{proof}
The Sugawara formula for $\widehat{\mathfrak g}_{k}$ gives
$c=k\dim\mathfrak g/(k+h^{\vee})$; at $k=1$ for
$\mathrm{SO}(10)$ this is $45/9=5$. The Abelian factor contributes $1$.
Their sum is $16$.
\end{proof}

\begin{theorem}[Heterotic worldsheet correspondence]
\label{thm:strings:het}
Combining the four right-moving spacetime bosons ($c_{R}=4$ Nambu--Goto)
with their $(0,1)$ worldsheet superpartners ($c_{R}=4\times\tfrac12=2$)
and the transverse fermions of the core gives a right-moving central
charge $c_{R}=10$; together with $c_{L}=16$ from \eqref{eq:strings:cL},
the vortex worldsheet realises the central-charge content
\begin{equation}
  (c_{L},c_{R})=(16,10)
  \label{eq:strings:cc}
\end{equation}
of the level-one left-moving sector of the $E_{8}\times E_{8}$ heterotic
string.
\end{theorem}

\begin{remark}[Correspondence, not duality]
\label{rem:strings:corr}
\autoref{thm:strings:het} is a \emph{low-energy correspondence}: the
gauge sector, anomaly structure, and central charges of the UCFT vortex
match the heterotic worldsheet at leading order in $\alpha'$. It is
\emph{not} a strict duality --- no full equivalence of S-matrices,
spectra, or quantum moduli spaces is claimed, and the matching is stated
only at the level of \eqref{eq:strings:cc} and the affine algebra
\eqref{eq:strings:affine}. The conformal-field-theory derivation
(Sugawara construction, modular invariance, higher-derivative
robustness) is deferred to \autoref{app:ws:main}.
\end{remark}

%----------------------------------------------------------------------%
\subsection{The $\mathbb{Z}_{5}$ quintic Calabi--Yau correspondence}
\label{sec:strings:cy}

The broken phase leaves $32 = 4+28$ real flat directions: the $4$
gauge/soldering clock fields $\pi_{-}^{i}$ ($i=0,1,2,3$) and the $28$
massive clock fields $\pi_{+}^{j}$ promoted to Wilson-line moduli once
the composite connection condenses,
\begin{equation}
  \mathcal{M}^{\mathrm{mod}}_{\mathrm{UCFT}}
   = \underbrace{\mathbb{R}^{4}/\Gamma^{4}}_{\text{Goldstone/soldering}}
     \times
     \underbrace{\mathbb{R}^{28}/\Gamma^{28}}_{\text{Wilson lines}},
   \qquad
   \dim_{\mathbb{R}}\mathcal{M}^{\mathrm{mod}}_{\mathrm{UCFT}}=32.
  \label{eq:strings:moduliU}
\end{equation}

On the heterotic side, compactify the $E_{8}\times E_{8}$ string on the
freely-acting $\mathbb{Z}_{5}$ quotient of the Fermat quintic,
$\widehat{Q}=Q/\mathbb{Z}_{5}$ with
$Q=\{\sum_{a=0}^{4}z_{a}^{5}=0\}\subset\mathbb{CP}^{4}$, standard
embedding of $T\widehat{Q}$ in one $E_{8}$:
\begin{equation}
  h^{1,1}(\widehat{Q})=1,\quad
  h^{2,1}(\widehat{Q})=0,\quad
  \pi_{1}(\widehat{Q})=\mathbb{Z}_{5},
  \label{eq:strings:hodge}
\end{equation}
the Wilson lines along $\pi_{1}(\widehat{Q})=\mathbb{Z}_{5}$ breaking
$E_{6}\to\mathrm{Spin}(10)\times U(1)$. Bundle cohomology gives
\begin{equation}
  h^{1}(\widehat{Q},V_{\mathbf{16}})=3,\qquad
  h^{1}(\widehat{Q},V_{\mathbf{10}})=14,\qquad
  h^{1}(\widehat{Q},\mathbf{1})=1,
  \label{eq:strings:bundle}
\end{equation}
hence the real-moduli count
\begin{equation}
  \underbrace{4}_{\text{K\"ahler}+B_{2}\text{ axions}}
  \;+\;
  \underbrace{28}_{2\times h^{1}(V_{\mathbf{10}})}
  \;=\;32
  \;=\;\dim_{\mathbb{R}}\mathcal{M}^{\mathrm{mod}}_{\mathrm{Het}} .
  \label{eq:strings:count}
\end{equation}

\begin{theorem}[Moduli matching]
\label{thm:strings:moduli}
The real dimensions and gauge-bundle topology of
$\mathcal{M}^{\mathrm{mod}}_{\mathrm{UCFT}}$ in
\eqref{eq:strings:moduliU} and of
$\mathcal{M}^{\mathrm{mod}}_{\mathrm{Het}}$ in \eqref{eq:strings:count}
agree: both are $32 = 4+28$, both carry an
$\mathrm{Spin}(10)\times U(1)$ bundle, and the kinetic metrics coincide
at tree level (each descends from the $E_{6}$ Killing form). The four
Goldstone/soldering coordinates map to the universal K\"ahler modulus
plus three axionic partners; the $28$ Wilson-line coordinates map to the
$2\times 14$ real bundle moduli in $H^{1}(\widehat{Q},V_{\mathbf{10}})$.
\end{theorem}

\begin{proof}[Proof sketch]
Counting and bundle topology are read off from
\eqref{eq:strings:moduliU}--\eqref{eq:strings:count}. Metric equality
follows from the common $E_{6}$ Killing origin of both kinetic terms;
Chern-class/Chern--Simons matching uses the shared Green--Schwarz
two-form. The detailed cohomology and the explicit map are given in
\autoref{app:cy:main}.
\end{proof}

\begin{remark}[Three generations: corroboration, not derivation]
\label{rem:strings:gen}
The chiral generation count of the heterotic model is
$h^{1}(\widehat{Q},V_{\mathbf{16}})=3$ from \eqref{eq:strings:bundle}.
This \emph{corroborates} the three generations of UCFT, which are
supplied as an added postulate $P5$ --- the horizontal $SU(3)_{F}$
symmetry on $J_{3}(\mathbb{O})\otimes\mathbb{C}^{3}$
(\autoref{sec:sm:main}) --- and \emph{not} as a second, independent
mechanism. The $14$ vector-like $\mathbf{10}$ moduli and the single
$\mathbf{1}$ (a right-handed-neutrino modulus) complete the
$32=4+28$ match. Generations are \emph{added structure} ($P5$), with the
heterotic count as a \emph{correspondence}.
\end{remark}

%----------------------------------------------------------------------%
\subsection{Conclusions: theorems versus postulates}
\label{sec:strings:conclusions}

UCFT is an $E_{6}$/Albert-algebra--motivated unification. Two axioms
supply the arena (the $27=J_{3}(\mathbb{O})$ order parameter and its
quantum kinematics); a small, explicitly enumerated set of postulates
$P1$--$P6$ then yields a four-dimensional Lorentzian quantum field theory
containing an anomaly-free $\mathrm{Spin}(10)\times U(1)$ gauge sector,
Sakharov-induced gravity, and --- after a stated GUT-Higgs/flavour sector
--- the Standard Model with three generations and neutrino masses. This
is a \emph{conditional embedding} of the Standard Model and gravitation.
\autoref{tab:strings:ledger} records the classification of results.

\begin{table}[h]
\centering
\renewcommand{\arraystretch}{1.25}
\begin{tabular}{lll}
\hline
\textbf{Result} & \textbf{Content} & \textbf{Status}\\
\hline
Vacuum selection ($P1\Rightarrow P2$)
  & global min $=$ rank-$1$ idempotent orbit $\mathrm{EIII}$, $\dim 32$
  & \textbf{Theorem}\\
Origin instability
  & $\mathrm{Hess}\,V(0)\prec0$; roll-down to $\mathcal{M}$
  & Theorem\\
Coset tangent / metric
  & $\mathfrak m=\mathbf{16}_{-3}\oplus\overline{\mathbf{16}}_{+3}$, $K|_{\mathfrak m}\succ0$
  & Theorem\\
Soldering algebra
  & $H_{2}(\mathbb{O})\cong\mathbb{R}^{1,9}$, $SL(2,\mathbb{O})=\mathrm{Spin}(1,9)$
  & Theorem\\
Signature $(1,3)$, $d=4$
  & added soldering structure
  & \textbf{Postulate $P3$}\\
Mass spectrum
  & $28$ massive (common $m^{2}=\mu^{2}$) $+\,4$ eaten
  & Theorem\\
Induced gravity
  & $M_{\mathrm{Pl}}^{2}=\tfrac{28}{6}f^{2}\mu^{2}\log(\Lambda^{2}/\mu^{2})/16\pi^{2}>0$
  & Theorem\\
Gauge running
  & asymptotically free, $b_{0}(1\,\mathrm{gen})=26$, $b_{0}(3\,\mathrm{gen})=12$
  & Theorem\\
Gravitational fixed point
  & $\kappa_{\ast}=96\pi^{2}/5$, $\mathrm{eig}=-2$
  & \textbf{Indicative}\\
SM descent
  & $\mathbf{16}=$ one family $+\,\nu^{c}$; anomalies cancel (GS)
  & Theorem\\
Chirality / generations / Yukawa
  & $P4,P5,P6$; seesaw $m_{\nu}\sim0.05$\,eV
  & \textbf{Postulate}\\
BPS vortex strings
  & finite tension $\mathcal{T}_{n}=2\pi f^{2}|n|\,\Omega$
  & Theorem\\
String topology
  & from gauged $U(1)$ config space, not $\pi_{1}(\mathrm{EIII})=0$
  & Theorem $+$ \textbf{Postulate $P$}\\
Heterotic worldsheet
  & $(c_{L},c_{R})=(16,10)$, $\widehat{\mathfrak{so}}(10)_{1}\oplus\widehat{\mathfrak u}(1)_{1}$
  & \textbf{Correspondence}\\
Calabi--Yau moduli
  & $32=4+28$, $h^{1}(V_{\mathbf{16}})=3$
  & \textbf{Correspondence}\\
Cosmological constant
  & tuned relevant $\Lambda_{cc}$
  & inherited, not solved\\
\hline
\end{tabular}
\caption{Classification of results by epistemic status. Theorems are
proved from stated hypotheses; postulates ($P1$--$P6$) are added
structure; conjectural results are motivated but outside controlled
weak coupling; correspondences are leading-order matches, not
dualities.}
\label{tab:strings:ledger}
\end{table}

\paragraph{What is proved.}
The vacuum-selection theorem ($P1\Rightarrow P2$, the rank-$1$
idempotent orbit $\mathcal{M}=\mathrm{EIII}$ of real dimension $32$), the
positive-definiteness and uniqueness of the coset metric, the octonionic
soldering algebra $SL(2,\mathbb{O})\cong\mathrm{Spin}(1,9)$, the common
radiative mass of the $28$ clock scalars by Schur's lemma, the positivity
$M_{\mathrm{Pl}}^{2}>0$ of the Sakharov-induced Planck mass with
prefactor $28/6$, the asymptotic freedom of the gauge--Yukawa sector
($\beta_{\hat g^{2}}=2\hat g^{2}+(b_{0}/48\pi^{2})\hat g^{4}$,
$b_{0}=12$ for three generations), the Green--Schwarz cancellation of the
mixed anomaly, the embedding of one Standard-Model family plus a
right-handed neutrino in the $\mathbf{16}$, and the finiteness of the BPS
vortex tension are all \emph{theorems}.

\paragraph{What is postulated.}
The Lorentzian signature and dimension ($P3$), the chiral matter content
without mirrors ($P4$), the three generations via $SU(3)_{F}$ ($P5$), and
the real GUT-Higgs/Yukawa sector generating the type-I seesaw ($P6$) are
\emph{added postulates}. The topological string sector rests on the
gauged-$U(1)$ configuration space (\autoref{post:strings:sector}), since
$\pi_{1}(\mathrm{EIII})=0$.

\paragraph{What is indicative or a correspondence.}
The gravitational fixed point $\kappa_{\ast}=96\pi^{2}/5\approx189.50$ is
\emph{indicative}: $\kappa_{\ast}\gg1$ lies outside controlled weak
coupling. Reflection positivity of the UCFT $\sigma$-model (and hence
Osterwalder--Schrader reconstruction of Lorentzian dynamics) is
indicative. The heterotic worldsheet match $(c_{L},c_{R})=(16,10)$ and
the $\mathbb{Z}_{5}$ quintic moduli match $32=4+28$ with
$h^{1}(V_{\mathbf{16}})=3$ are \emph{low-energy correspondences}, not
strict dualities. UCFT \emph{inherits, and does not solve}, the
cosmological-constant problem: $\Lambda_{cc}$ is a tuned relevant
coupling.

\paragraph{Outlook.}
Several directions would sharpen the embedding into a derivation. (i) A
curved-background functional-RG truncation could test whether the
indicative gravitational fixed point survives beyond the single-scalar
loop, and whether the higher-curvature operators are genuinely
irrelevant. (ii) A direct check of reflection positivity for the gauged
clock-field $\sigma$-model would upgrade emergent unitary time from
indicative to theorem. (iii) Deriving $d=4$ and signature $(1,3)$ from a
dynamical vierbein-condensate selection, rather than postulating $P3$,
remains open. (iv) A worldsheet computation of the full
vortex partition function would test the heterotic correspondence beyond
central charges. (v) Phenomenologically, dimension-six proton decay
$\tau(p\to e^{+}\pi^{0})\sim M_{\mathrm{GUT}}^{4}/(\alpha_{\mathrm{GUT}}^{2}
m_{p}^{5})$, leptogenesis from the heavy $\nu^{c}$, and inflation driven
by one of the $28$ massive clock fields after the GUT transition all
furnish falsifiable tests. UCFT is presented as a conditional embedding:
the proved core---one Albert-algebra order parameter selecting an
anomaly-free $\mathrm{Spin}(10)\times U(1)$ gauge sector with induced,
positive Planck mass---is established under the stated hypotheses, while the
Lorentzian, flavour, and gravitational-UV completions rest on explicit
postulates and conjectures enumerated above.
